\chapter{The same transforms in exact arithmetic}

\textbf{One transform here is demonstrated; the other two are design only, and the predictions are
unmeasured.} The translation by number theoretic transform is built and run in
\texttt{examples/\allowbreak{}0\_experimental/\allowbreak{}exact\_translation\_by\_ntt.py}: its
correlation agrees to the digit with the direct \(O(N^2)\) count and recovers the planted shift. Its
grade is \emph{agrees}, a run against a reference and not a prediction. The precision constants it rests
on are re-derived and a wrong one refused in
\texttt{examples/\allowbreak{}0\_experimental/\allowbreak{}ntt\_twiddle\_certificate.py}. The
integer-field and rotation transforms are not built: nothing in their sections has been measured, and
they are the design this work adds to the methods quoted before it, stated before any result. The
predictions at the end of the chapter concern behavior on real sequences and remain unmeasured for
every transform, the translation included.

\section{Why exactness changes the order of operations}

Every published method in this book rounds. A fast Fourier transform carries each coefficient to
machine precision, a sub-pixel peak is fitted by a curve, a variational problem is solved to a
tolerance, and a spherical integral is a quadrature on a mesh. Each rounding is small, and each is
fed to the next stage: the translation that was rounded is removed before the rotation is estimated,
and the rotation that was rounded is removed before the flow is. A stratified pipeline is exactly the
arrangement in which small errors compound, because every stage reads the residual the last stage
left, including what the last stage got wrong.

The consequence for design is specific. Averaging exists to cancel error: fitting a line across many
frames, pooling neighbors, smoothing a field. Where the arithmetic makes no error there is nothing to
cancel, and averaging only mixes the present with the past and one object with another. What remains
useful is accumulation: exact counts added over time, which lose nothing however many are added.

\section{Integer fields}

A smoothing kernel with integer weights keeps a field in the integers. The binomial kernel of even
order \(n\), the \(n\)th row of Pascal's triangle, sums to exactly \(2^n\) and has variance \(n/4\) in units of
the grid spacing. It is the integer counterpart of a Gaussian of that variance. Separable along each
axis, a volume smoothed by orders \((n_1, \dots, n_d)\) is scaled by exactly \(2^{n_1 + \dots + n_d}\), and
the difference of two smoothings,
\[
r \;=\; s \cdot 2^{\,\sum_i m_i} \;-\; b,
\]
with \(s\) smoothed by orders \(n_i\) and \(b\) smoothed further by orders \(m_i\), compares them at one common
scale with nothing divided. Its width is fixed by the orders and the input depth. It is carried in a
fixed number of limbs, and its sign says exactly whether a point stands above its local background.

\section{Translation by number theoretic transforms}

The agreement at every lag, \(C(\ell) = \sum_v a(v)\, b(v + \ell)\), is a convolution of the reflected
first view with the second. A number theoretic transform computes a convolution exactly in the
integers modulo a prime \(p\) that carries roots of unity of the transform length. The prime is
\(p = 119 \cdot 2^{23} + 1 = 998244353\), with \(p - 1 = 2^{23} \cdot 7 \cdot 17\). 3 is a primitive
root and the longest power-of-two transform the modulus admits has length \(2^{23}\), past any image
axis. The certificate is re-derived in
\texttt{examples/\allowbreak{}0\_experimental/\allowbreak{}ntt\_twiddle\_certificate.py}, and the prime,
its generator and the \(2^{27}\) device primes beside it are proved in
\texttt{theory/\allowbreak{}theory/\allowbreak{}twiddle\_constants\_article.tex}. This prime is below \(2^{31}\),
which the same proof shows is the property that keeps a butterfly's sum and its subtraction from
wrapping in a 32-bit word; a modulus above \(2^{31}\) carried a silent wrong answer there. A single
small prime keeps this transform clear of that boundary.

Exactness has a stated bound. Each coefficient is a sum of products \(a(v)\, b(v + \ell)\). For binary
views it counts agreeing positions and is at most the transform length. The transform length divides
\(p - 1\) and so is below \(p\) for any valid prime, which makes a binary view exact with no separate
precondition: no residue is ever reduced. Weighting the views is what can breach \(p\). The bound becomes
the sum of the weight products, and past \(p\) a single prime counts modulo itself, silently, the same
fault one level up from a wrong constant. The remedy is a prime above the largest coefficient, or the
Chinese remainder theorem over several, the device path of that proof reassembling a 94-bit product. A
view is embedded on each axis in a period of at least twice its extent less one so that no two lags
alias, and the transform is separable, one axis at a time. The construction holds in any number of
dimensions.

The peak is then an integer, and ties among lags are broken by a rule, not by rounding: the shortest
lag by a weighted squared length whose weights carry the view's anisotropy, then a fixed index order.

\section{Rotation without angles}

A rotation in the plane is multiplication by a unit complex number, and the ratio of two directions
is a complex number whose argument is the turn between them. Where directions are given by rational
coordinates the ratio is a Gaussian rational, and applying it is exact. A turn carried forward one
step is exact in the same way. What is not rational is a fraction of a turn, a cube root of a
rotation for instance, and a design that needs one has left exact arithmetic and should say so.

For turns in three dimensions the same holds for quaternions with rational components. The spherical
harmonic route of the third chapter needs coefficients that are integrals, and a function sampled at
points gives those integrals only through a quadrature rule; which quadratures are exact for which
degrees on which point sets is the question an exact design has to answer first.

\section{Predictions}

\textbf{First.} Removing the view's motion before linking will reduce links made to the wrong object
more than any change to the linker does, wherever the view moves by as much as the objects do,
because in that regime most of every object's apparent step belongs to the view.

\textbf{Second.} A view that orbits its subject leaves a residual after the best translation that is
structured, not scattered: parallax that varies with depth, or a turn that varies with distance from
an axis. The residual of the translation stage is the measurement that says the next transform is
needed.

\textbf{Third.} Each transform recovered exactly and removed will leave the next stage a residual with
no rounding in it. The improvement from adding a stage will not shrink the way it does in a floating
point pipeline, where each stage inherits the error of the last.
